% Consolidated guide-ch14 from f8f4709b0ffac9da4c1fecbf2ce873002f6da6cc; source fragments preserved.
% 2026-11-24 八批严格二校第5批：历史片段哈希对应修改前版本。

% === BEGIN guide-ch14-examples.tex blob f8838d9986d7ffe6066a72fd5c60ee475b553d9c ===
% 第13批：《大学物理学学习指导（第5版）》第14章“光的偏振”。
% 例14.1—14.4；印刷199—201页／PDF物理207—209页。
% 依源图及数字独立推导，不复制原书整段文字和受保护原图。
\subsection{例题：马吕斯定律、混合偏振光与布儒斯特角}

\sourceexample{14.1}{199}{207}
\problemtext 自然光依次通过两个透振轴夹角为$60^\circ$的理想偏振片，
原透射强度为$I_1$。在这两个偏振片之间再插入一片，
它与两侧透振轴均成$30^\circ$，求新透射强度$I$。
\workings
设自然光入射强度为$I_0$。通过第一片后强度为$I_0/2$，
再经过原第二片，由马吕斯定律
\[
 I_1=\frac{I_0}{2}\cos^260^\circ=\frac{I_0}{8},
 \qquad I_0=8I_1.
\]
加入中间的偏振片后，两次相邻轴夹角均为$30^\circ$，
\[
 I=\frac{I_0}{2}\cos^230^\circ\cos^230^\circ
 =\frac{I_0}{2}\left(\frac34\right)^2
 =\boxed{\frac94 I_1=2.25I_1}.
\]
\crosscheck 新强度虽高于原来两片的透射强度，
仍小于入射自然光的$I_0/2$，
符合对偏振方向逐步旋转所产生的透射变化。

\sourceexample{14.2}{199--200}{207--208}
\problemtext 两片偏振片$P_1,P_2$叠放，照射光由强度相等的
自然光和完全线偏振光混合而成。
通过$P_1$后的强度为总入射强度的一半，
通过两片后的强度为总入射强度的四分之一。
求入射线偏振成分与$P_1$轴夹角$\theta$和两片轴夹角$\alpha$；
再计及每片对理应透过分量均额外吸收$5\%$，求新的角度。
\workings
令两种入射光分量的强度各为$I_0$，
因此总入射强度$I_{\rm in}=2I_0$。
不计额外吸收时，第一片透射
\[
 I_{P_1}=\frac12I_0+I_0\cos^2\theta
 =\frac12I_{\rm in}=I_0.
\]
于是$\cos^2\theta=\tfrac12$，即
$\boxed{\theta=45^\circ}$。
第二片后的强度
$I_{P_2}=I_{P_1}\cos^2\alpha
 =I_{\rm in}/4=I_0/2$，
故$\boxed{\alpha=45^\circ}$。

若每片对其理应透过的光另吸收$5\%$，
保持原先给定的\emph{两个测得强度占入射总强度的比值}，
须有
\[
 0.95 I_0\left(\frac12+\cos^2\theta\right)=I_0,
 \quad
 \cos^2\theta=\frac1{0.95}-\frac12=\frac{21}{38},
\]
\[
 \boxed{\theta=\arccos\sqrt{\frac{21}{38}}
 \approx42.0^\circ}.
\]
第二片的吸收还给出
\[
 0.95\cos^2\alpha=\frac12,\quad
 \cos^2\alpha=\frac{10}{19},\qquad
 \boxed{\alpha=\arccos\sqrt{\frac{10}{19}}
 \approx43.5^\circ}.
\]
\crosscheck 首片同时接受自然光与线偏振光，不能仅以
$I_{\rm in}\cos^2\theta$处理整个混合入射束。
第二片两次透射的比值为$1/2$，因此吸收模型中
须对每片分别记入$0.95$，不可把两次损耗重复施加于同一片。

\sourceexample{14.3}{200--201}{208--209}
\problemtext 三种透明介质$I,II,III$有互相平行的界面，
折射率$n_1=1.00$、$n_2=1.43$、$n_3$待求。
自然光从$I$入射到$II$，又由$II$入射到$III$；
两个界面的反射光均达到完全线偏振。
求入射角$i$和$n_3$。
\workings
第一界面达到布儒斯特条件，
\[
 \tan i=\frac{n_2}{n_1}=1.43,\qquad
 \boxed{i=\arctan1.43\approx55.03^\circ}.
\]
其折射角$\gamma$满足$i+\gamma=90^\circ$，
且因界面平行，$\gamma$又是第二界面的入射角。
第二界面亦处在布儒斯特角，
\[
 \tan\gamma=\frac{n_3}{n_2}
 \quad\Longrightarrow\quad
 n_3=n_2\cot i
 =\frac{n_2n_1}{n_2}=\boxed{1.00}.
\]
\crosscheck 两次完全偏振都要求$p$分量反射系数为零。
对不吸收、非磁、各向同性介质，该条件成立，
与源图中的双平行界面排列一致。

\sourceexample{14.4}{201}{209}
\problemtext 一平板玻璃置于水中，与水面夹角为$\theta$。
空气、水、玻璃折射率分别为$1,1.333,1.517$；
要求水面及玻璃面产生的两束反射光均完全线偏振，
求夹角$\theta$（按原图光线在两面间的几何路径）。
\workings
空气到水界面的布儒斯特角为
$i_1=\arctan(1.333)\approx53.12^\circ$。
水到玻璃界面的布儒斯特角为
\[
 i_2=\arctan\frac{1.517}{1.333}\approx48.69^\circ.
\]
第一界面的折射角
$\gamma=90^\circ-i_1$。
由源图中两界面的几何三角形关系，
$\theta=i_2-\gamma$，因此
\[
 \boxed{\theta=i_1+i_2-90^\circ\approx11.8^\circ}.
\]
若未限定原图法线方位，另一共面几何分支给锐角$\gamma+i_2\approx85.6^\circ$，需按实际光路判定可达性。
\crosscheck 此式取源图所示光线从空气入水、再在
倾斜玻璃面反射的方位分支。
若玻璃板倾角取反或换入射路径，应重新绘制法线并确定
有向角；不能仅凭两个布儒斯特角相减而忽略图中几何。
% === END guide-ch14-examples.tex ===
